\subsection{Datasets and Experimental Conditions}
\label{sec:datasets}

To evaluate CESM-Diff under zero-shot scenarios, we conduct experiments on three widely adopted industrial benchmarks: Tennessee Eastman Process (TEP), Hydraulic System, and CWRU Bearing datasets. These benchmarks collectively span failure modes from continuous chemical processes to hydraulic subsystem degradations and rotating machinery faults, providing a comprehensive testbed for zero-shot fault diagnosis~\cite{lei2020applications, zhao2019deep}. A summary of dataset statistics is provided in Table~\ref{tab:dataset_stats}.

\subsubsection{Benchmark Datasets}

\textbf{Tennessee Eastman Process (TEP).} The TEP dataset, originally proposed by Downs and Vogel~\cite{downs1993plant}, simulates a realistic chemical plant comprising five major unit operations: a reactor, a condenser, a vapor--liquid separator, a recycle compressor, and a product stripper. The simulation generates 52 process variables (22 continuous measurements and 19 composition measurements sampled every 3 minutes, along with 11 manipulated variables) under one normal operating condition and 21 programmed fault types. Each fault type is characterized by a distinct root cause, including step changes in feed composition, random variations in cooling water temperature, slow drift in reaction kinetics, and sticking valves, covering a spectrum from abrupt to incipient faults with varying severity levels. Following the partition strategy used in prior zero-shot fault diagnosis studies~\cite{feng2021fault}, we select 15 fault classes and organize them into five experimental groups. In each group, 12 classes serve as seen faults for training (480 samples per class), and 3 classes are held out as unseen faults for zero-shot evaluation. The non-linear cross-correlations among variables and significant process noise in TEP pose substantial challenges for generative models, which CESM-Diff addresses through its CLIP-based semantic anchoring mechanism.

\textbf{Hydraulic System.} This dataset is collected from a real-world hydraulic test rig instrumented with 17 physical sensors that monitor pressure (6 sensors at 100~Hz), motor power (1 sensor at 100~Hz), volume flow (2 sensors at 10~Hz), temperature (4 sensors at 1~Hz), vibration (1 sensor at 1~Hz), virtual cooling efficiency (1 sensor at 1~Hz), virtual cooling power (1 sensor at 1~Hz), and system efficiency factor (1 sensor at 1~Hz). The rig comprises four primary subsystems---cooler, valve, pump, and accumulator---each of which can exhibit multiple degradation levels. Health states are quantitatively defined: the cooler condition ranges from full efficiency (100\%) to close to total failure (3\%); the valve condition ranges from optimal switching (100\%) to severe lag (73\%); the pump leakage is graded from no leakage to severe internal leakage (2~mL/min); and the accumulator pressure spans from optimal (130~bar) to critically low (90~bar). By combining these subsystem states, we define 12 distinct health conditions, with 9 as seen classes and 3 as unseen classes. A total of 2205 measurement cycles are recorded. The multi-scale nature of the sensor signals---spanning sampling rates from 1~Hz to 100~Hz---requires CESM-Diff to align physically heterogeneous sensor modalities within a unified CLIP-derived semantic manifold, testing the model's capacity for cross-rate feature fusion.

\textbf{CWRU Bearing.} The Case Western Reserve University (CWRU) bearing dataset~\cite{cwru_dataset} is a cornerstone benchmark in mechanical fault diagnosis. Vibration signals are collected from a 2-hp Reliance Electric motor driving a shaft equipped with SKF deep-groove ball bearings. Accelerometers mounted at the drive end record vibration data at a sampling frequency of 12~kHz under four shaft speeds (1730, 1750, 1772, and 1797~RPM) corresponding to motor loads of 3, 2, 1, and 0~hp, respectively. Single-point faults are seeded at three locations---inner raceway, outer raceway, and rolling element---using electro-discharge machining, with fault diameters of 0.007, 0.014, and 0.021 inches representing minor, moderate, and severe damage. Raw 1-D vibration signals are segmented using a sliding window of 1024 points and subsequently transformed into 2-D time--frequency representations via Short-Time Fourier Transform (STFT) to capture jointly the temporal and spectral characteristics of fault signatures. We define 10 unique fault classes (including the normal condition): 7 serve as seen classes and 3 as unseen classes. This configuration tests the model's ability to distinguish fine-grained mechanical failure modes using only textual semantic embeddings generated by the CLIP encoder.

\subsubsection{Experimental Strategy}
All experiments are implemented in PyTorch on an NVIDIA RTX 3090 GPU. Input signals are min-max normalized to $[-1, 1]$ based on statistics computed exclusively from seen-class training data. Sliding windows of 1024 points are applied for CWRU and Hydraulic datasets, while TEP uses the native 3-minute sampling intervals. For diffusion, $T=1000$ steps with a linear noise schedule are employed during training. The U-Net backbone adopts a 4-layer 1-D residual architecture with channel dimensions $[64, 128, 256, 512]$. The CLIP text encoder (ViT-B/32) remains frozen throughout training. Optimization uses AdamW with an initial learning rate of $5 \times 10^{-4}$, cosine annealing schedule, and a batch size of 256. Strict data separation is enforced: no unseen-class samples appear during training or hyperparameter tuning. Each experiment is repeated 10 times with different random seeds, and we report Mean Class Accuracy (MCA) with standard deviations.
